\documentclass[11pt]{article}
\usepackage[margin=1in]{geometry}
\usepackage{amsmath,amssymb,amsthm,booktabs,array}
\usepackage[T1]{fontenc}
\usepackage{lmodern}
\usepackage[hidelinks]{hyperref}
\hypersetup{pdftitle={Paper V: Corrections and Scoped Lean Completion},pdfauthor={Komorebi P.},pdfsubject={Checked scoped formalization, statement corrections, and verification boundaries}}
\usepackage{fancyhdr}
\pagestyle{fancy}
\fancyhf{}
\lhead{Projector-Resolved Seed Calculus V}
\rhead{Komorebi P.}
\cfoot{\thepage}
\setlength{\headheight}{14pt}
\newtheorem{proposition}{Proposition}
\newcommand{\C}{\mathbb C}
\newcommand{\R}{\mathbb R}
\title{Paper V\\Corrections and Scoped Lean Completion}
\author{Komorebi P.}
\date{Version 1.0.0, 6 September 2026}
\begin{document}
\maketitle
\noindent\textbf{Status.} This supplement accompanies the preserved 31-page adaptive-tree-resolved manuscript. The corrected, scope-explicit formalization contract now covers all 20 theorem-like results, using 225 Lean theorem targets in 23 modules. The original 133 checked targets remain included; 92 targets have been added, including 11 foundations adapted from Paper IV. These counts include supporting lemmas and are not counts of new manuscript theorems.

The accompanying verification report records the build and axiom audit. The allowed axioms are \texttt{propext}, \texttt{Classical.choice}, and \texttt{Quot.sound}.

Correspondence between the manuscript and the formal statements is manually reviewed; Lean checks the formal statements, not the PDF. Read the manuscript together with the corrections and scope boundaries below.

\section{Finite selection, including a singleton tree}
Theorem 11.3, page 19, requires a separate case when no competing pruning exists. A root-only tree has precisely one complete pruning, so a minimum taken over all competing prunings is not a real number in that case.

\begin{proposition}[Corrected finite selection statement]
Let $\mathcal P$ be a nonempty finite set of complete prunings, and let $J:\mathcal P\to\R$. A global minimizer exists. If $\mathcal P=\{\Lambda_*\}$, then $\Lambda_*$ is the unique minimizer of every real-valued objective on $\mathcal P$.

Suppose instead that $|\mathcal P|\ge2$ and $\Lambda_*$ uniquely minimizes $J$. Then
\[
\gamma=\min_{\Lambda\in\mathcal P\setminus\{\Lambda_*\}}
\bigl(J(\Lambda)-J(\Lambda_*)\bigr)>0.
\]
For another objective $J'$, if
\[
|J'(\Lambda)-J(\Lambda)|\le\varepsilon
\quad\hbox{for every }\Lambda\in\mathcal P,
\qquad 2\varepsilon<\gamma,
\]
then $\Lambda_*$ also uniquely minimizes $J'$.
\end{proposition}
\begin{proof}
A real-valued function on a nonempty finite set attains its minimum. The singleton case has no competitor. In the nonsingleton case, uniqueness makes every difference in the nonempty finite competitor set positive, so its minimum is positive. For each competitor,
\begin{align*}
J'(\Lambda)-J'(\Lambda_*)
&\ge J(\Lambda)-J(\Lambda_*)-2\varepsilon\\
&\ge\gamma-2\varepsilon>0.
\end{align*}
\end{proof}
The inequality must be strict. The objectives $(0,2)$ and $(1,1)$ have uniform difference one, exactly half the original gap, and the second objective has a tie.

\section{A deterministic tree selector and its envelope}
For Theorem 11.4, pages 19--20, specify the tie rule:
\[
\text{stop at }v\quad\Longleftrightarrow\quad
c_v\le\sum_{u\in\operatorname{ch}(v)}V(u).
\]
Terminal nodes are always selected. Strictly smaller child cost causes refinement. Bottom-up induction shows that each selected subtree cost is its minimum. With the costs already supplied, one addition per child and one comparison per internal node give a linear arithmetic-operation count. Computing the local costs and the bit complexity of exact arithmetic are separate costs.

For local errors $|c'_v-c_v|\le\varepsilon_v$, define
\[
E(v)=
\begin{cases}
\varepsilon_v,&v\text{ terminal},\\
\max\left\{\varepsilon_v,\sum_{u\in\operatorname{ch}(v)}E(u)\right\},&v\text{ internal}.
\end{cases}
\]
Induction over any complete pruning gives
\[
\left|\sum_{v\in\Lambda}c'_v-\sum_{v\in\Lambda}c_v\right|\le E(\mathrm{root}).
\]
The finite-gap result applies when $2E(\mathrm{root})<\gamma$. In Corollary 11.5, the probability statement should refer to a \emph{joint event} on which all the local bounds hold. Separate per-node probability statements are not automatically a joint bound with the same failure probability.

Exact synthesis is independent of the selected objective. Repeated substitution of $P_v=\sum_{u\in\operatorname{ch}(v)}P_u$ stops at a complete pruning and yields $\sum_{v\in\Lambda}P_v=I$. The checked construction covers arbitrary positive finite branching, with a two-way correspondence between recursive cuts and set-valued covering antichains. The synthesis identity is a literal finite sum of selected node operators.

\section{Channel delay and original-sample delay}
In Theorem 4.1, page 8, let $\zeta$ denote the \emph{channel} transform variable. With
\[
x_r[q]=x[Nq+r],
\]
the identity $R(\zeta)E(\zeta)=\zeta^{-d}I_N$ yields
\[
\widehat x_r[q]=x_r[q-d],\qquad
\widehat x[n]=x[n-Nd].
\]
It therefore describes $d$ channel-index steps, or $Nd$ original samples.

For a general original-sample delay $D=q_0N+t$, where $0\le t<N$, put
\[
s_r=(r-t)\bmod N,\qquad d_r=q_0+\mathbf1_{\{r<t\}}.
\]
Then the exact delay matrix has the entries
\[
\mathcal D_D(\zeta)_{r,s}=
\begin{cases}
\zeta^{-d_r},&s=s_r,\\
0,&s\ne s_r.
\end{cases}
\]
Indeed, $Nq+r-D=N(q-d_r)+s_r$. For two channels and one original sample,
\[
\mathcal D_1(\zeta)=
\begin{pmatrix}0&\zeta^{-1}\\1&0\end{pmatrix},
\]
not $\zeta^{-1}I_2$. Thus either name the delay unit in the existing scalar criterion or use $R E=\mathcal D_D$ for arbitrary original-sample delays.

The time-domain FIR law does not need an analytic transform. Its transform restatement should use finite Laurent polynomials, another specified formal interpretation, or an actual common convergence domain.

\section{The Smith action at zero and infinity}
The finite rational formula in Theorem 13.2, page 23, is not an everywhere-defined map $\C\to\C$. In particular, the half-turn sends zero impedance to infinity.

\begin{proposition}[Projective phase action]
For $w\in\C\setminus\{0\}$, let
\[
M(w)=\begin{pmatrix}1+w&1-w\\1-w&1+w\end{pmatrix}.
\]
Its projective action on $\mathbb P^1(\C)$ is invertible, and these actions satisfy $T_u\circ T_v=T_{uv}$. For $w=-1$, the action is $[a:b]\mapsto[b:a]$.
\end{proposition}
\begin{proof}
Direct multiplication gives
\[
\det M(w)=4w\ne0,\qquad M(u)M(v)=2M(uv).
\]
The nonzero scalar factor does not change a projective map. Also
$M(-1)=\left(\begin{smallmatrix}0&2\\2&0\end{smallmatrix}\right)$.
\end{proof}
For $w=\omega_N^k$, this gives the cyclic composition law. On the affine chart $[z:1]$ where the denominator is nonzero, it gives exactly the printed formula. Corollary 13.3 becomes $z\mapsto z^{-1}$ for $z\ne0$, with the projective extension exchanging zero and infinity.

Alternatively, the finite action can be restricted to $\operatorname{Re}z>0$. Since
\[
|z+1|^2-|z-1|^2=4\operatorname{Re}z>0,
\]
its Cayley image lies strictly inside the unit disk, which is preserved by phase rotations. The finite closed passive half-plane omits the disk point $\Gamma=1$, corresponding to infinite impedance. Ordinary division by zero in a formal library must not be used to encode this projective point.

\section{Other hypothesis and editorial repairs}
For Theorem 7.1, specify measurable masks and membership of the relevant indicator functions in $L^p$. For finite $p$, finite measures of the original mask and its symmetric difference suffice for the displayed finite-norm formula. Distinguish a pointwise supremum from an essential supremum.

On page 28, change the implementation table's sixteen figure pairs to eighteen after checking the companion source. On page 30, Appendix A.3, change ``Harmically fill'' to ``Harmonically fill.'' These corrections do not promote either edge-recovery conjecture to a theorem.

\section{Transform interpretation and exact filter-bank scope}
The time-domain identity in Theorem 3.1 holds for every bilateral sequence $x$ and every finitely supported filter. To read its z-transform identity as an equality of complex numbers, fix $z\ne0$ and require each channel series
\[
X_s(z)=\sum_{n\in\mathbb Z}x_s[n]z^{-n}
\]
to be summable. The checked theorem then proves summability of the output channel series and
\[
Y_k(z)=\sum_{s=0}^{N-1}\left(\sum_m h[Nm+k-s]z^{-m}\right)X_s(z).
\]
The inner sum is finite. No convergence condition is imposed on the time-domain identity, and no divergent series is assigned a transform by convention.

For the concrete operator form of Theorem 4.1, finite Laurent polynomials represent FIR channel filters. The monomial $z^{-d}$ acts by $x[n]\mapsto x[n-d]$. The formalization proves that this action is faithful: an impulse recovers every coefficient. It also proves that multiplication of filter-bank matrices corresponds exactly to composition of the time-domain banks. Consequently, reconstruction of every bilateral channel sequence with gain $c$ and channel delay $d$ is equivalent to
\[
R E=c z^{-d}I_N.
\]
This concrete FIR statement needs no transform convergence. Where transforms converge, evaluating the same Laurent matrices gives the channel-transform formula. An arbitrary scalar-matrix annihilation criterion is also checked; the release does not assert convergence or boundedness for unspecified infinite filters.

\section{What closes the remaining formalization gaps}
The Haar energy statement now uses the actual samples $x[2n]$ and $x[2n+1]$. Their channel energies sum to the original $\ell^2$ norm. At every finite depth, the constructed Haar packet operators satisfy both the full energy identity and synthesis using their actual Hilbert-space adjoints.

The adaptive-tree results now use literal finite sums over set-valued covering antichains. The recursive cuts and those antichains are equivalent. The nonadditive objective has a minimizer; a unique minimizer with competitors has an attained positive gap. A cached bottom-up evaluator returns the global optimum, stops on ties, and uses exactly one comparison per internal node.

For the complexity statement in Theorem 11.4, scalar node costs are supplied as inputs. Charging one unit for a node visit, a comparison, and an addition of a cached child value gives a proved bound of $3|T|$ for these operations. This is a specified arithmetic-operation model, not a bound on Lean runtime, arbitrary cost-oracle computation, or printing the selected tree.

The perturbation theorem is checked for measurable masks on arbitrary measure spaces using extended $L^p$ norms. Finite measures of the original mask and its symmetric difference give the displayed finite-valued statement. Both the pointwise supremum bound and its essential-supremum consequence are proved. The noisy-selector probability clause follows from inclusion of the joint local-error event.

The nonlinear cyclic formula is checked for formal power series over a commutative semiring, including the literal $q$-tuple power formula and $q=0$. Analytic evaluations retain their ordinary convergence requirements. The Smith action uses the explicit exponential root of unity and a quotient of nonzero homogeneous pairs by nonzero scaling, so its modulo-$N$ law is an actual projective action. The $N=2$ specialization is connected to $e^{\pi i}=-1$.

\section{Verification and what remains outside the contract}
The release verifies 225 named theorem targets, all imported project-namespace declarations, the pinned Lean and dependency revisions, source integrity, and principal-target coverage for all 20 contract rows. The independent mathematical suites retain 20 passing tests with 305,933 counted exact assertions; 16 verifier controls check failure handling, including incomplete scope rows and unregistered targets. These finite tests do not replace the universal Lean proofs.

Across Papers I--V, the census finds 1,443 mapped proof declarations after removing two shared Paper I/II theorem names and the 11 adapted Paper IV proofs. There are also 319 distinct mapped constructions in Papers I--IV; Paper V's additional definitions are not included in that construction count. The five manuscripts contain 168 theorem/lemma/proposition/corollary records. Supporting declarations and printed results are different counting units, and neither count establishes mathematical novelty.

The earlier Papers I--IV build status is taken from their bundled release audits; their sources were not freshly rebuilt in this census. Paper V was built and audited here. The manuscript's original 33-test companion suite, reported depth-five optimum, numerical residuals, and priority claims are outside this completion contract. The preserved manuscript has not been silently rewritten: this supplement supplies the authoritative corrections. The edge-recovery conjectures remain conjectures.
\end{document}
